\Appendix

\includegraphics[width=\textwidth]{appendix-all-1}

\includegraphics[width=\textwidth]{appendix-all-2}

\includegraphics[width=\textwidth]{appendix-all-3}

\includegraphics[width=0.9\textwidth]{appendix-all-4}
% \includegraphics[width=0.9\textwidth]{appendix-all-5}
% \begin{center}
%   \large
%   数学专业师范生教学技能调查问卷
% \end{center}

% 亲爱的同学，你们好！这是一份关于数学专业师范生教学技能的调查问卷，本次问卷调查大约需要2-3分钟时间，请大家认真耐心的填写!填写问卷的同学会在审核后收到2元红包，并致以谢意!您的回答会严格保密，只用于研究，不会泄露给任何人，感谢您的支持！

% 第一部分：基本信息

% \begin{enumerate}
%   \item 性别：①男生  ②女生
%   \item 年级：①大三  ②大四
%   \item 计算机等级：①一级  ②二级  ③三级及以上
%   \item 普通话等级：①一甲  ②一乙  ③二甲  ④二乙
% \end{enumerate}

% 第二部分：

% 一、教学设计技能

% 5. 请根据您的实际情况选择最符合的项：1-->5表示“不具备-->完全具备”

% \includegraphics[width = \textwidth]{appendix-5.png}
% % \begin{tblr}{
% %   width = 0.98\textwidth,
% %   colspec = {X[9.8,l]X[1,c]X[1,c]X[1,c]X[1,c]X[1,c]},
% %   hlines, vlines,
% %   row{odd} = {bg=azure9},
% %   row{1} = {bg=azure8}
% % }
% % & 1 & 2 & 3 & 4 & 5 \\
% % ①能够制定合理、恰当的教学目标 & ○ & ○ & ○ & ○ & ○ \\
% % ②能够深入钻研、梳理和全面分析教材 & ○ & ○ & ○ & ○ & ○ \\
% % ③能够准确分析、预测学生情况 & ○ & ○ & ○ & ○ & ○ \\
% % ④能够选择合适的教学手段方法 & ○ & ○ & ○ & ○ & ○ \\
% % ⑤能够预设教学过程 & ○ & ○ & ○ & ○ & ○ 
% % \end{tblr}

% 二、数学说课技能

% 6. 请根据您的实际情况选择最符合的项：1-->5表示“不具备->完全具备”

% \includegraphics[width = \textwidth]{appendix-6.png}
% % \begin{tblr}{
% %   width = 0.98\textwidth,
% %   colspec = {X[9.8,l]X[1,c]X[1,c]X[1,c]X[1,c]X[1,c]},
% %   hlines, vlines,
% %   row{odd} = {bg=azure9},
% %   row{1} = {bg=azure8}
% % }
% % & 1 & 2 & 3 & 4 & 5 \\
% % ①清楚怎样说课，了解说课程序 & ○ & ○ & ○ & ○ & ○ \\
% % ②说课紧扣教学大纲 & ○ & ○ & ○ & ○ & ○ \\
% % ③说课时语言简明、表达清晰、重点突出 & ○ & ○ & ○ & ○ & ○ \\
% % {④能够系统而概括的传达自己的教学观点、\\教学方法、理论依据等} & ○ & ○ & ○ & ○ & ○ \\
% % \end{tblr}


% 三、课堂教学技能

% 7. 语言技能 

% 【请根据您的实际情况选择最符合的项：1-->5表示“不具备-->完全具备”】

% \includegraphics[width = \textwidth]{appendix-7.png}

% % \begin{tblr}{
% %   width = \textwidth,
% %   colspec = {X[11,l]X[1,c]X[1,c]X[1,c]X[1,c]X[1,c]},
% %   hlines, vlines,
% %   row{odd} = {bg=azure9},
% %   row{1} = {bg=azure8}
% % }
% % & 1 & 2 & 3 & 4 & 5 \\
% % ①音量大小、语速、节奏合适 & ○ & ○ & ○ & ○ & ○ \\
% % ②语言目的明确，紧扣学习内容，建立知识联系 & ○ & ○ & ○ & ○ & ○ \\
% % ③语言处理灵活，教学效果好，具有指导性 & ○ & ○ & ○ & ○ & ○ \\
% % \end{tblr}

% 8.板书技能 

% 【请根据您的实际情况选择最符合的项：1-->5表示“不具备-->完全具备”】

% \includegraphics[width = \textwidth]{appendix-8.png}

% % \begin{tblr}{
% %   width = 0.97\textwidth,
% %   colspec = {X[9.8,l]X[1,c]X[1,c]X[1,c]X[1,c]X[1,c]},
% %   hlines, vlines,
% %   row{odd} = {bg=azure9},
% %   row{1} = {bg=azure8}
% % }
% % & 1 & 2 & 3 & 4 & 5 \\
% % ①板书简洁、概括、直观性强 & ○ & ○ & ○ & ○ & ○ \\
% % ②板书字迹清晰、大小合适、重点突出 & ○ & ○ & ○ & ○ & ○ \\
% % ③能科学的体现知识结构及基本思路和方法 & ○ & ○ & ○ & ○ & ○ \\
% % \end{tblr}


% 9.讲解技能

% 【请根据您的实际情况选择最符合的项：1-->5表示“不具备-->完全具备”】

% \includegraphics[width = \textwidth]{appendix-9.png}

% % \begin{tblr}{
% %   width = 0.98\textwidth,
% %   colspec = {X[9.8,l]X[1,c]X[1,c]X[1,c]X[1,c]X[1,c]},
% %   hlines, vlines,
% %   row{odd} = {bg=azure9},
% %   row{1} = {bg=azure8}
% % }
% % & 1 & 2 & 3 & 4 & 5 \\
% % ①讲解内容丰富，符合学生认知水平 & ○ & ○ & ○ & ○ & ○ \\
% % ②讲解言简意赅、通俗易懂，富有启发性 & ○ & ○ & ○ & ○ & ○ \\
% % ③讲解思路清晰，重点突出 & ○ & ○ & ○ & ○ & ○ \\
% % ④讲解能够及时反馈、强化、巩固学生所学 & ○ & ○ & ○ & ○ & ○ \\
% % \end{tblr}


% 10.提问技能

% 【请根据您的实际情况选择最符合的项：1-->5表示“不具备-->完全具备”】

% \includegraphics[width = \textwidth]{appendix-10.png}

% % \begin{tblr}{
% %   width = \textwidth,
% %   colspec = {X[10,l]X[1,c]X[1,c]X[1,c]X[1,c]X[1,c]},
% %   hlines, vlines,
% %   row{odd} = {bg=azure9},
% %   row{1} = {bg=azure8}
% % }
% % & 1 & 2 & 3 & 4 & 5 \\
% % ①提问后注意停顿，让学生有必要的思考时间 & ○ & ○ & ○ & ○ & ○ \\
% % ②所提问题明确、清晰、有启发性、难易程度合适 & ○ & ○ & ○ & ○ & ○ \\
% % ③提问面广，能够照顾到各类学生 & ○ & ○ & ○ & ○ & ○ \\
% % \end{tblr}


% 11.导入技能

% 【请根据您的实际情况选择最符合的项：1-->5表示“不具备-->完全具备”】

% \includegraphics[width = \textwidth]{appendix-11.png}

% % \begin{tblr}{
% %   width = 0.98\textwidth,
% %   colspec = {X[9.8,l]X[1,c]X[1,c]X[1,c]X[1,c]X[1,c]},
% %   hlines, vlines,
% %   row{odd} = {bg=azure9},
% %   row{1} = {bg=azure8}
% % }
% % & 1 & 2 & 3 & 4 & 5 \\
% % ①导入目的明确，方法与新知识联系紧密 & ○ & ○ & ○ & ○ & ○ \\
% % ②导入能引起学生的学习兴趣，集中注意力 & ○ & ○ & ○ & ○ & ○ \\
% % ③能够根据具体情况采取不同的导入方法 & ○ & ○ & ○ & ○ & ○ \\
% % ④能够较好的把握导入时间 & ○ & ○ & ○ & ○ & ○ \\
% % \end{tblr}


% 12.结束技能

% 【请根据您的实际情况选择最符合的项：1-->5表示“不具备-->完全具备”】

% \includegraphics[width = \textwidth]{appendix-12.png}

% % \begin{tblr}{
% %   width = 0.98\textwidth,
% %   colspec = {X[9.8,l]X[1,c]X[1,c]X[1,c]X[1,c]X[1,c]},
% %   hlines, vlines,
% %   row{odd} = {bg=azure9},
% %   row{1} = {bg=azure8}
% % }
% % & 1 & 2 & 3 & 4 & 5 \\
% % ①结束方式与内容相适应 & ○ & ○ & ○ & ○ & ○ \\
% % ②结束目的明确，使知识条理化、系统化 & ○ & ○ & ○ & ○ & ○ \\
% % ③结束能够检查学习，强化学习效果 & ○ & ○ & ○ & ○ & ○ \\
% % \end{tblr}


% 13.媒体运用技能

% 【请根据您的实际情况选择最符合的项：1-->5表示“不具备-->完全具备”】

% \includegraphics[width = \textwidth]{appendix-13.png}

% % \begin{tblr}{
% %   width = 0.98\textwidth,
% %   colspec = {X[9.8,l]X[1,c]X[1,c]X[1,c]X[1,c]X[1,c]},
% %   hlines, vlines,
% %   row{odd} = {bg=azure9},
% %   row{1} = {bg=azure8}
% % }
% % & 1 & 2 & 3 & 4 & 5 \\
% % ①能够熟练制作课件、教学视频等 & ○ & ○ & ○ & ○ & ○ \\
% % ②能够熟练使用数学教学软件（几何画板、Geogebra等） & ○ & ○ & ○ & ○ & ○ \\
% % ③能够熟练寻找教学资源 & ○ & ○ & ○ & ○ & ○ \\
% % ④能够根据教学内容、学生情况等因素选择合适的媒体进行教学 & ○ & ○ & ○ & ○ & ○ \\
% % \end{tblr}


% 四、教学评价技能

% 14. 请根据您的实际情况选择最符合的项：1-->5表示“不具备-->完全具备”

% \includegraphics[width = \textwidth]{appendix-14.png}

% % \begin{tblr}{
% %   width = 0.98\textwidth,
% %   colspec = {X[9.8,l]X[1,c]X[1,c]X[1,c]X[1,c]X[1,c]},
% %   hlines, vlines,
% %   row{odd} = {bg=azure9},
% %   row{1} = {bg=azure8}
% % }
% % & 1 & 2 & 3 & 4 & 5 \\
% % ①掌握评课基本方法，能够分析、归纳、总结一堂课的优缺点，并提出改进建议 & ○ & ○ & ○ & ○ & ○ \\
% % ②能够从多个维度和不同角度进行评课 & ○ & ○ & ○ & ○ & ○ \\
% % ③能够采取适用于教学决策的评价方法 & ○ & ○ & ○ & ○ & ○ \\
% % \end{tblr}


% 五、教学研究技能

% 15. 请根据您的实际情况选择最符合的项：1-->5表示“不具备-->完全具备”

% \includegraphics[width = \textwidth]{appendix-15.png}

% % \begin{tblr}{
% %   width = 0.98\textwidth,
% %   colspec = {X[9.8,l]X[1,c]X[1,c]X[1,c]X[1,c]X[1,c]},
% %   hlines, vlines,
% %   row{odd} = {bg=azure9},
% %   row{1} = {bg=azure8}
% % }
% % & 1 & 2 & 3 & 4 & 5 \\
% % ①掌握数学研究基本方法，了解研究基本过程 & ○ & ○ & ○ & ○ & ○ \\
% % ②具备查找、整理、分析与概括文献的能力 & ○ & ○ & ○ & ○ & ○ \\
% % ③能够了解学科发展动态和吸取科研成果 & ○ & ○ & ○ & ○ & ○ \\
% % \end{tblr}


% 六、教学技能培养情况

% 16. 你认为下列活动培养了哪方面教学技能？（多选）

% \includegraphics[width = \textwidth]{appendix-16.png}

% \begin{tblr}{
%   width = 0.98\textwidth,
%   colspec = {X[9.8,c]X[1,c]X[1,c]X[1,c]X[1,c]X[1,c]},
%   hlines, vlines,
%   row{odd} = {bg=azure9},
%   row{1} = {bg=azure8}
% }
% &教学设计技能	&数学说课技能	&课堂教学技能	&教学评价技能	&教学研究技能\\
% 教学比赛、师范生技能比赛 & □ & □ & □ & □ & □ \\
% 教育见习和实习 & □ & □ & □ & □ & □ \\
% 微格教学 & □ & □ & □ & □ & □ \\
% 理论学习（教育学、心理学、数学教学论） & □ & □ & □ & □ & □ \\
% 观摩学习（录像、优秀教师课堂） & □ & □ & □ & □ & □ \\
% \end{tblr}


% 17. 到目前为止，参与下列活动的次数。
\hspace*{1em}
\includegraphics[width = 0.8\textwidth]{appendix-17.png}

% \begin{tblr}{
%   width = 0.98\textwidth,
%   colspec = {X[5,l]X[1,c]X[1,c]X[1,c]X[1,c]X[1,c]},
%   hlines, vlines,
%   row{odd} = {bg=azure9},
%   row{1} = {bg=azure8}
% }
% &0次	&0-5次&	5-10次	&10次以上\\
% 教学设计 & ○ & ○ & ○ & ○ \\
% 说课 & ○ & ○ & ○ & ○ \\
% 讲课 & ○ & ○ & ○ & ○ \\
% 评课 & ○ & ○ & ○ & ○ \\
% 教学研究 & ○ & ○ & ○ & ○ \\
% \end{tblr}

